\documentclass[conference]{IEEEtran}
\usepackage[utf8]{inputenc}
\usepackage[T1]{fontenc}
\usepackage{graphicx}
\usepackage{booktabs}
\usepackage{amsmath}
\usepackage{url}
\usepackage{xcolor}
\usepackage{caption}
\usepackage{hyperref}
\hypersetup{colorlinks=true, linkcolor=black, citecolor=blue!60!black, urlcolor=blue!60!black}

% numbers and table fragments produced by scripts/07_report_assets.py
\IfFileExists{generated/numbers.tex}{% generated by scripts/07_report_assets.py - do not edit
\newcommand{\maseNaiveHOne}{0.471}
\newcommand{\smapeNaiveHOne}{13.3}
\newcommand{\maseNaiveHSix}{2.094}
\newcommand{\smapeNaiveHSix}{54.5}
\newcommand{\maseNaiveHDay}{0.873}
\newcommand{\smapeNaiveHDay}{24.2}
\newcommand{\maseNaiveAll}{1.146}
\newcommand{\smapeNaiveAll}{30.7}
\newcommand{\maseSeasonalNaiveHOne}{0.854}
\newcommand{\smapeSeasonalNaiveHOne}{22.3}
\newcommand{\maseSeasonalNaiveHSix}{0.863}
\newcommand{\smapeSeasonalNaiveHSix}{23.0}
\newcommand{\maseSeasonalNaiveHDay}{0.873}
\newcommand{\smapeSeasonalNaiveHDay}{24.2}
\newcommand{\maseSeasonalNaiveAll}{0.863}
\newcommand{\smapeSeasonalNaiveAll}{23.2}
\newcommand{\maseWeeklyNaiveHOne}{2.380}
\newcommand{\smapeWeeklyNaiveHOne}{51.0}
\newcommand{\maseWeeklyNaiveHSix}{2.369}
\newcommand{\smapeWeeklyNaiveHSix}{51.0}
\newcommand{\maseWeeklyNaiveHDay}{2.326}
\newcommand{\smapeWeeklyNaiveHDay}{51.5}
\newcommand{\maseWeeklyNaiveAll}{2.358}
\newcommand{\smapeWeeklyNaiveAll}{51.2}
\newcommand{\maseSarimaHOne}{0.399}
\newcommand{\smapeSarimaHOne}{11.0}
\newcommand{\maseSarimaHSix}{1.542}
\newcommand{\smapeSarimaHSix}{35.9}
\newcommand{\maseSarimaHDay}{2.341}
\newcommand{\smapeSarimaHDay}{54.5}
\newcommand{\maseSarimaAll}{1.427}
\newcommand{\smapeSarimaAll}{33.8}
\newcommand{\maseLightgbmHOne}{0.949}
\newcommand{\smapeLightgbmHOne}{26.1}
\newcommand{\maseLightgbmHSix}{1.780}
\newcommand{\smapeLightgbmHSix}{43.0}
\newcommand{\maseLightgbmHDay}{1.568}
\newcommand{\smapeLightgbmHDay}{40.6}
\newcommand{\maseLightgbmAll}{1.432}
\newcommand{\smapeLightgbmAll}{36.6}
\newcommand{\maseGruHOne}{1.012}
\newcommand{\smapeGruHOne}{26.8}
\newcommand{\maseGruHSix}{2.699}
\newcommand{\smapeGruHSix}{60.0}
\newcommand{\maseGruHDay}{2.957}
\newcommand{\smapeGruHDay}{62.9}
\newcommand{\maseGruAll}{2.223}
\newcommand{\smapeGruAll}{49.9}
\newcommand{\bestModel}{SARIMAX}
\newcommand{\bestModelMase}{1.427}
\newcommand{\bestOverallModel}{Seasonal naive (24 h)}
\newcommand{\bestOverallMase}{0.863}
\newcommand{\naiveDayMB}{275}
\newcommand{\naiveAllGB}{16.7}
\newcommand{\rawDayMB}{308}
\newcommand{\rawRowsM}{4.5}
\newcommand{\ingestPeakMB}{1,551}
\newcommand{\ingestTotalMin}{5.0}
\newcommand{\ingestDays}{62}
\newcommand{\pipelinePeakMB}{5,244}
\newcommand{\stlDailyShare}{81}
\newcommand{\stlResidShare}{6}
\newcommand{\mstlDailyShare}{81}
\newcommand{\mstlWeeklyShare}{9}
\newcommand{\giniInternet}{0.62}
\newcommand{\topTenShare}{50}
\newcommand{\topOneShare}{12}
\newcommand{\internetShare}{82}
\newcommand{\nHours}{1488}
\newcommand{\selectedCells}{duomo (\#5060), bocconi (\#4259), navigli (\#4456), business (\#4855), business-2 (\#5161), mixed suburban (\#5857), mixed suburban-2 (\#4755)}
\newcommand{\nSelected}{7}
\newcommand{\bestK}{2}
\newcommand{\bestSilhouette}{0.47}
\newcommand{\chosenK}{4}
\newcommand{\chosenSilhouette}{0.28}
\newcommand{\optunaTrials}{20}
\newcommand{\optunaBestMase}{0.633}
\newcommand{\figdir}{../reports/figures}
}{\newcommand{\figdir}{../reports/figures}}
\providecommand{\bestModel}{(pending)}
\providecommand{\bestModelMase}{--}
\providecommand{\bestOverallModel}{(pending)}
\providecommand{\bestOverallMase}{--}
\providecommand{\selectedCells}{(pending)}
\providecommand{\nSelected}{--}
\providecommand{\giniInternet}{--}
\providecommand{\topTenShare}{--}
\providecommand{\topOneShare}{--}
\providecommand{\internetShare}{--}
\providecommand{\nHours}{--}
\providecommand{\bestK}{--}
\providecommand{\bestSilhouette}{--}
\providecommand{\chosenK}{--}
\providecommand{\chosenSilhouette}{--}
\providecommand{\naiveDayMB}{--}
\providecommand{\naiveAllGB}{--}
\providecommand{\rawDayMB}{--}
\providecommand{\rawRowsM}{--}
\providecommand{\ingestPeakMB}{--}
\providecommand{\ingestTotalMin}{--}
\providecommand{\ingestDays}{--}
\providecommand{\pipelinePeakMB}{--}
\providecommand{\stlDailyShare}{--}
\providecommand{\stlResidShare}{--}
\providecommand{\mstlDailyShare}{--}
\providecommand{\mstlWeeklyShare}{--}
\providecommand{\optunaTrials}{--}
\providecommand{\optunaBestMase}{--}

% Passages that must be re-read against the real-data run before submission.
% Set \revisemode to 0 to print them as normal text.
\newcommand{\revisemode}{0}
\newcommand{\revise}[1]{\ifnum\revisemode=1 {\color{red!70!black}#1}\else #1\fi}

\newcommand{\fig}[1]{\figdir/#1}

\begin{document}
\raggedbottom

\title{Forecasting Mobile Network Traffic in Milan:\\ Does Land Use Decide Which Model Wins?}

\author{\IEEEauthorblockN{Samuel Ihimbazwe}
\IEEEauthorblockA{African Leadership University\\
s.ihimbazwe@alustudent.com\\
Code: \url{https://github.com/Samkwizera/Time-Series-Forecasting}}}

\maketitle

\begin{abstract}
Short-term forecasts of cellular traffic help operators allocate capacity and switch off idle equipment, but each cell follows the activity pattern of the area it serves. This study compares three forecasting families: seasonal ARIMA (SARIMAX), gradient-boosted trees (LightGBM), and recurrent neural networks (LSTM/GRU). We use the 20.8\,GB Telecom Italia Milan dataset, covering 10\,000 grid cells at 10-minute resolution for 62 days. The pipeline streams and compresses the raw files, with a measured peak RSS of \pipelinePeakMB{}\,MB, and evaluates the models on \nSelected{} representative cells at 1, 6 and 24-hour horizons. The test period covers Christmas and New Year. The 24-hour seasonal naive was best overall (mean MASE \bestOverallMase). Among the learned models, \bestModel{} was best (\bestModelMase), narrowly ahead of LightGBM, and performance varied more by horizon than by the descriptive cell labels.
\end{abstract}

\begin{IEEEkeywords}
mobile traffic forecasting, time series, SARIMA, gradient boosting, recurrent neural networks, Telecom Italia Big Data Challenge
\end{IEEEkeywords}

% ------------------------------------------------------------------------------------
\section{Introduction}
Cellular networks are dimensioned for peak demand, yet demand in any given cell swings by an order of magnitude over a day. If the load of the coming hours can be predicted, an operator can pre-allocate spectrum, steer users between overlapping cells, or switch off carriers overnight to save energy. Errors have asymmetric costs: an under-forecast at the evening peak means dropped sessions, while an over-forecast at 4\,a.m.\ wastes energy. The forecasting problem is therefore practical, but it is also heterogeneous: a cell covering an office block, a residential street or a nightlife district exhibits a different daily shape, and a single citywide model may not serve all of them equally well.

The Telecom Italia Big Data Challenge dataset \cite{telecomitalia2015,barlacchi2015} makes this question tractable. It records SMS, call and Internet activity for a 100$\times$100 grid over Milan at 10-minute resolution for two months. At 20.8\,GB it is also large enough that the first engineering decision -- how to load it at all on a machine with 12--30\,GB of RAM -- is part of the problem.

\textbf{Objective.} To build a reproducible pipeline for hourly Internet-activity forecasting on this dataset and to compare three model families with different inductive biases under a common protocol.

\textbf{Research question.} \emph{How does forecasting accuracy differ between a statistical model (SARIMAX), a gradient-boosted tree model (LightGBM) and a recurrent network (LSTM/GRU) across grid cells with different activity-profile labels and across 1\,h, 6\,h and 24\,h horizons?} The cell-level comparison distinguishes this study from earlier work on the same data, which typically reports one aggregate error over all cells. The labels are inferred from traffic profiles and are not verified land-use categories.

% ------------------------------------------------------------------------------------
\section{Related Work}
Barlacchi et al.\ \cite{barlacchi2015} describe the dataset and note its strong daily and weekly periodicity and the concentration of activity in the city centre. Furno et al.\ \cite{furno2017} show, across ten cities, that mobile-traffic ``signatures'' -- the normalised weekly profile of a cell -- cluster into a small number of land-use classes (residential, office, transport, leisure). This motivates our clustering step: instead of hand-picking landmarks, we let the profile shape define cell types and then ask whether the forecasting ranking changes across them.

Most forecasting papers on the Milan data use deep spatio-temporal models. Wang et al.\ \cite{wang2017} combine autoencoders with LSTMs to capture spatial dependence between cells; Zhang et al.\ \cite{zhang2018} treat the grid as an image and apply densely connected CNNs; Huang et al.\ \cite{huang2017} compare RNN, 3D-CNN and CNN-RNN variants. These works report that deep models beat ARIMA baselines, but the baselines are usually untuned non-seasonal ARIMA and comparisons are made on aggregate metrics. Makridakis et al.\ \cite{makridakis2018} warn that, in controlled comparisons, well-specified statistical methods often match or beat machine-learning methods on short univariate series, and Januschowski et al.\ \cite{januschowski2022} document why gradient-boosted trees have won recent forecasting competitions: they handle calendar covariates and heterogeneous series with little tuning. Our design follows these cautions -- the statistical baseline is a properly seasonal SARIMAX, the tree model receives the same calendar information as the network, and every model is judged against seasonal-naive forecasts with scale-free MASE \cite{hyndman2006} and Diebold--Mariano tests \cite{diebold1995,harvey1997}, as recommended by Hyndman and Athanasopoulos \cite{hyndman2021}.

% ------------------------------------------------------------------------------------
\section{Dataset and Data Preparation}
\subsection{Data}
The dataset \cite{telecomitalia2015} consists of 62 daily tab-separated files (1 Nov 2013 -- 1 Jan 2014), 20.8\,GB in total, with columns \texttt{square\_id}, \texttt{time\_interval} (Unix ms, UTC), \texttt{country\_code} and five activity measures (\texttt{sms\_in}, \texttt{sms\_out}, \texttt{call\_in}, \texttt{call\_out}, \texttt{internet}). Values are proportional to CDR counts, rescaled by an undisclosed constant. Each (cell, interval) pair appears once per country code, so a day holds roughly \rawRowsM{} million rows (\rawDayMB{}\,MB).

\subsection{Memory management}
Loading a single day with default \texttt{pandas} dtypes needs an estimated \naiveDayMB{}\,MB; the full period would need about \naiveAllGB{}\,GB, several times the RAM of a Colab or Kaggle session. Three decisions bring this down (Table~\ref{tab:memory}):
\begin{enumerate}
\item \textbf{Stream and aggregate early.} Each file is scanned lazily with Polars' streaming engine (a chunked \texttt{pandas} fallback is provided), and rows are summed over \texttt{country\_code} before anything is materialised. The country dimension is irrelevant to load forecasting and removing it drops roughly 90\,\% of the rows.
\item \textbf{Downcast.} Activities become \texttt{float32} and cell ids \texttt{uint16}; timestamps are converted to Europe/Rome local time once, at ingestion, so that ``8\,a.m.'' means the same thing in every downstream plot.
\item \textbf{Persist compact intermediates.} One zstd-compressed Parquet file per day is written, and from these a wide hourly matrix (time $\times$ 10\,000 cells, \texttt{float32}) is built per activity. Only these matrices are read by the analysis and modelling stages; the raw TSVs are touched exactly once.
\end{enumerate}
Peak resident memory during the per-day ingestion loop was \ingestPeakMB{}\,MB and the \ingestDays{} days took \ingestTotalMin{} minutes. The later citywide 10-minute aggregation reached the pipeline maximum of \pipelinePeakMB{}\,MB. Both figures fit within the Kaggle environment, but the second shows that early aggregation did not keep every stage below 2\,GB. Every stage records peak RSS and wall time in \texttt{reports/tables/memory\_log.csv}; the same log supplies the compute-cost comparison in Section~\ref{sec:results}.

\begin{table}[t]
\centering\caption{Memory footprint of the ingestion strategy (RSS = resident set size).}\label{tab:memory}
\footnotesize
\IfFileExists{generated/memory_table.tex}{\begin{tabular}{p{4.6cm} rr}
\toprule
Stage & RSS (MB) & Time (s) \\
\midrule
Naive pandas load, one day (est.) & 275 & -- \\
Naive pandas load, all 62 days (est.) & 17,058 & -- \\
Streaming aggregation, per day (median) & 1,420 & 4.9 \\
Streaming aggregation, all days (max, total) & 1,551 & 300 \\
wide hourly internet & 1,514 & 17.5 \\
citywide 10min & 5,244 & 9.5 \\
eda & 567 & 10.9 \\
tsa & 554 & 56.3 \\
\bottomrule
\end{tabular}
}{(run scripts/07\_report\_assets.py)}
\end{table}

\subsection{Preprocessing for forecasting}
Missing (cell, interval) pairs denote zero activity and are filled with 0. Daily files are cut at UTC midnight, so the same local hour can appear in two files; the hourly matrix is therefore re-aggregated by timestamp after concatenation. The forecasting target is hourly Internet activity, which accounts for \internetShare{}\,\% of all recorded activity and is the quantity most relevant to capacity planning.

% ------------------------------------------------------------------------------------
\section{Exploratory Analysis}
\subsection{Temporal structure}
Fig.~\ref{fig:citywide} shows citywide hourly totals over the full period. The dominant feature is the daily cycle, modulated by a weekly cycle with lower weekend activity; the last week of December departs from both. Fig.~\ref{fig:profiles} decomposes the mean daily profile by day type. \revise{On weekdays Internet activity rises steeply from 6\,a.m., plateaus through office hours and peaks in the evening; weekends start later and lack the morning ramp; public holidays resemble Sundays. SMS and calls peak earlier and fall away in the evening, when Internet activity is still high -- evidence that the activities are driven by different behaviours and should not be pooled into one target.} The hour-by-weekday heatmap (Fig.~\ref{fig:hourdow}) makes the weekly rhythm explicit and shows that Friday and Saturday nights extend later than other nights.

\begin{figure}[t]\centering
\includegraphics[width=\columnwidth]{\fig{eda_citywide_hourly}}
\caption{Citywide hourly activity for all five measures. Shaded bands are Italian public holidays (1 Nov, 7--8 Dec, 25--26 Dec, 1 Jan).}\label{fig:citywide}
\end{figure}
\begin{figure}[t]\centering
\includegraphics[width=\columnwidth]{\fig{eda_daily_profiles}}
\caption{Mean hourly profile by day type for each activity.}\label{fig:profiles}
\end{figure}
\begin{figure}[t]\centering
\includegraphics[width=\columnwidth]{\fig{eda_hour_dow_internet}}
\caption{Mean citywide Internet activity by hour of day and day of week.}\label{fig:hourdow}
\end{figure}

\subsection{Anomalous days}
Daily totals with a robust $z$-score computed against the same weekday (Fig.~\ref{fig:daily}) flag the days that a model trained on ordinary weeks will find hardest. \revise{The flagged days coincide with public holidays -- 1 November, Sant'Ambrogio (7 December, a Milan-only holiday), 8 December and the Christmas period -- confirming that a holiday indicator is informative and that a test window covering Christmas is a stress test rather than a representative sample.}

\begin{figure}[t]\centering
\includegraphics[width=\columnwidth]{\fig{eda_daily_totals_internet}}
\caption{Daily Internet totals; labelled bars have $|z|>3.5$ relative to the same weekday.}\label{fig:daily}
\end{figure}

\subsection{Spatial structure}
Fig.~\ref{fig:spatial} maps mean Internet activity and weekday snapshots at four hours. Activity is extremely concentrated: the busiest 1\,\% of cells carry \topOneShare{}\,\% of Internet activity and the busiest 10\,\% carry \topTenShare{}\,\% (Gini coefficient \giniInternet). \revise{The centre (Duomo, the central station and the business district around Porta Nuova) dominates during office hours, whereas at 9\,p.m.\ the ring of residential districts lights up relative to the core. The suburban and rural cells in the outer grid show almost no activity at any hour.} This heterogeneity is the reason to evaluate per cell type rather than on the citywide aggregate, whose smoothness makes any model look good.

\begin{figure}[t]\centering
\includegraphics[width=\columnwidth]{\fig{eda_spatial_internet}}
\caption{Spatial distribution of Internet activity (log scale): overall weekday mean and snapshots at 04:00, 11:00, 15:00 and 21:00.}\label{fig:spatial}
\end{figure}

\subsection{Cell types from profile shape}
Following \cite{furno2017}, each cell above the median activity level is represented by its normalised weekday 24-hour profile and clustered with $k$-means. The silhouette criterion favours $k=\bestK$ (\bestSilhouette). We retain $k=\chosenK$ (silhouette \chosenSilhouette) because the finer partition separates two business-shaped and two mixed/suburban-shaped profiles used for the cell-level comparison. Fig.~\ref{fig:clusters} shows these four centroids and their spatial layout. This is an interpretability choice rather than an optimum under the silhouette metric. Cluster names come from centroid shape, using peak hour and the day/evening ratio, and do not constitute external land-use labels.

For the forecasting experiments we select \nSelected{} cells: the landmark cells used in the dataset's companion material (Duomo, Bocconi University, Navigli) and the busiest cell of each cluster (\selectedCells). Their hourly series are shown in Fig.~\ref{fig:selected}.

\begin{figure}[t]\centering
\includegraphics[width=\columnwidth]{\fig{eda_clusters_internet}}
\caption{Cluster centroids of normalised weekday profiles (left) and their spatial layout (right).}\label{fig:clusters}
\end{figure}
\begin{figure}[t]\centering
\includegraphics[width=\columnwidth]{\fig{eda_selected_cells_internet}}
\caption{Hourly Internet activity of the selected cells.}\label{fig:selected}
\end{figure}

\subsection{Time-series properties}
Table~\ref{tab:stationarity} reports ADF \cite{dickey1979} and KPSS \cite{kwiatkowski1992} tests on the citywide series. \revise{The raw and log series are non-stationary in the KPSS sense (the level drifts with the seasonal cycle and the December decline), while a seasonal difference at lag 24 renders the series stationary under both tests; a further first difference is unnecessary.} This decides the SARIMA differencing orders ($d=0$, $D=1$) and motivates the $\log(1+y)$ transform used by all models, which stabilises the variance that grows with the level.

STL decomposition \cite{cleveland1990} with a 24-hour period attributes \stlDailyShare{}\,\% of the variance of the log series to the daily component and \stlResidShare{}\,\% to the remainder; MSTL \cite{bandara2021} with periods 24 and 168 assigns \mstlDailyShare{}\,\% to the daily and \mstlWeeklyShare{}\,\% to the weekly component (Fig.~\ref{fig:mstl}). The daily cycle is thus the first-order effect, the weekly cycle a second-order one, and the remainder small but concentrated around holidays.

The ACF (Fig.~\ref{fig:acf}) shows slowly decaying peaks at multiples of 24 and a pronounced peak at 168; after seasonal differencing the PACF has a negative spike at lag 24, the signature of a seasonal MA term, and significant low-order lags. These observations motivate the SARIMA candidates, the tree-model lag set (1--24, then daily steps from 36 to 168), and a one-week candidate input window for the recurrent models. Validation, rather than the plots alone, determines the final specification.

\begin{table}[t]
\centering\caption{Stationarity tests on the citywide hourly Internet series.}\label{tab:stationarity}
\scriptsize\setlength{\tabcolsep}{3pt}
\IfFileExists{generated/stationarity_table.tex}{\begin{tabular}{l rr rr l}
\toprule
Transform & ADF & $p$ & KPSS & $p$ & Verdict \\
\midrule
raw & -2.14 & 0.230 & 3.63 & 0.010 & non-stationary \\
log1p & -1.26 & 0.649 & 3.70 & 0.010 & non-stationary \\
log1p, diff 1 & -11.91 & 0.000 & 0.00 & 0.100 & stationary \\
log1p, seasonal diff 24 & -5.58 & 0.000 & 0.22 & 0.100 & stationary \\
log1p, diff 1 + seasonal diff 24 & -9.82 & 0.000 & 0.01 & 0.100 & stationary \\
\bottomrule
\end{tabular}
}{(pending)}
\end{table}
\begin{figure}[t]\centering
\includegraphics[width=\columnwidth]{\fig{tsa_mstl_citywide}}
\caption{MSTL decomposition of the log citywide series into trend, daily and weekly seasonal components and remainder.}\label{fig:mstl}
\end{figure}
\begin{figure}[t]\centering
\includegraphics[width=\columnwidth]{\fig{tsa_acf_citywide}}
\caption{ACF/PACF of the log citywide series before and after a 24-hour seasonal difference. Dotted lines at lags 24, 48, 72, 168.}\label{fig:acf}
\end{figure}

% ------------------------------------------------------------------------------------
\section{Methodology}
\subsection{Forecasting setup}
Given hourly Internet activity $y_t$ of a cell, a forecast origin $t$ and a horizon $h\in\{1,6,24\}$, a model predicts $\hat y_{t+h}$ from data up to $t$ (direct strategy). The split is chronological and identical for all models: training 1 Nov -- 14 Dec, validation 15 -- 21 Dec, test 22 Dec -- 1 Jan. The test window deliberately contains Christmas and New Year; forecasting the holiday fortnight from six ordinary weeks is the situation an operator faces every year and is where the models are expected to differ most. Forecasts are issued from every hour of the evaluation window (rolling origin), so each (model, cell, horizon) is scored on roughly 240 forecasts.

\subsection{Input representation and preprocessing}
All models operate on $z_t=\log(1+y_t)$, standardised per cell with training-set statistics for the tree and recurrent models; predictions are transformed back before scoring. Candidate specifications include hour, weekday, weekend and holiday information (Italian national holidays plus Sant'Ambrogio). LightGBM receives these as columns and the recurrent models use cyclical encodings at each time step. The SARIMAX search tested weekly Fourier terms and a holiday dummy, but validation selected the specification without either regressor.

\subsection{Models}
\textbf{Baselines.} Naive ($\hat y_{t+h}=y_t$), 24-hour seasonal naive and 168-hour seasonal naive. The 24-hour seasonal naive is also the scaling series for MASE.

\textbf{SARIMAX.} The search covers seasonal ARIMA models $(p,0,q)(P,1,Q)_{24}$ on $z_t$, fitted separately to each cell \cite{seabold2010}. Weekly Fourier regressors and a holiday dummy are optional candidates because a seasonal period of 168 would make the state vector impractically large. Validation selected $(1,0,1)(1,1,1)_{24}$ without exogenous regressors (run s2). Parameters are estimated once on the fitting window; the Kalman filter then processes the evaluation window with fixed parameters and supplies an $h$-step forecast from each origin.

\textbf{LightGBM.} One gradient-boosted regressor per horizon \cite{ke2017} is trained jointly on the selected cells with a categorical cell identifier and an $\ell_1$ objective. The search adds feature groups in stages: lags 1--24 and 36, 48, \dots, 168; origin and target-time calendar attributes; and rolling statistics and differences. Validation selected run g3, which uses the full lag set and calendar variables but excludes rolling features. Early stopping on the last week of the fitting window fixes the number of trees, after which the model is refit on the whole window. For the final test run, the fitting window is train plus validation.

\textbf{Recurrent network.} The recurrent search compares LSTM \cite{hochreiter1997} and GRU \cite{cho2014} cells in a shared PyTorch implementation \cite{paszke2019}. Each model reads 24 or 168 hours of $(z_t, \text{calendar}_t)$ and predicts all three horizons from the final hidden state, a learned cell embedding and the calendar encoding of each target time. Training uses Adam, $\ell_1$ loss, gradient clipping, a plateau learning-rate schedule and early stopping. Run l5, a one-layer GRU with 64 hidden units and 168 hours of context, achieved the best recurrent validation MASE and was selected for the test set.

\textbf{Why these three.} SARIMAX represents a linear seasonal process, LightGBM learns nonlinear interactions among lags and calendar variables, and the recurrent family learns a representation of the recent sequence. The first two reflect strong classical and tabular forecasting approaches \cite{januschowski2022}; recurrent models also allow comparison with prior work on this dataset \cite{wang2017,huang2017}. We exclude spatial CNNs because evaluating neighbour information fairly would require a different input and sampling design.

\subsection{Evaluation}
Errors are reported in original units as MAE, RMSE, sMAPE and MASE \cite{hyndman2006}, per cell and horizon. MASE $<1$ means the model beats the 24-hour seasonal naive forecast on the training scale, which makes cells of very different volume comparable. Statistical significance of differences in absolute error is assessed with the Diebold--Mariano test with the Harvey--Leybourne--Newbold small-sample correction and $h-1$ Newey--West lags \cite{diebold1995,harvey1997}. Training time and peak memory are recorded per model.

\subsection{Experimental design and tuning strategy}
Tuning follows an iterative protocol: each run is scored on the validation split, the result is written to an experiment log together with the reasoning for the next change, and the test split is touched only once per model with the configuration selected on validation. The SARIMAX and recurrent rounds are manual because each change tests a hypothesis about the data (Table~\ref{tab:log}). For LightGBM, manual rounds compare feature groups before \optunaTrials{} trials of TPE search \cite{akiba2019} tune capacity and regularisation.

\begin{table*}[t]
\centering\caption{Experiment log (validation split). The rationale column is written before each run; the full log with all parameters is in \texttt{experiments/experiment\_log.md}.}\label{tab:log}
\scriptsize
\IfFileExists{generated/experiment_log_table.tex}{\begin{tabular}{l l rr r p{10.4cm}}
\toprule
Model & Run & val MASE & val sMAPE & s & Rationale for the run \\
\midrule
SARIMAX & s1 & 0.717 & 12.1 & 3 & Minimal seasonal random walk + AR(1): the simplest structure consistent with the ACF (slow decay at multiples of 24). Reference point for what the seasonal difference alone buys. \\
SARIMAX & s2 & 0.659 & 10.8 & 44 & Add MA terms at lag 1 and 24: PACF after seasonal differencing shows a negative spike at 24, the classic signature of a seasonal MA(1). \\
SARIMAX & s3 & 0.666 & 11.1 & 108 & Weekly cycle cannot be modelled with a 168-period SARIMA (state too large). Add 3 Fourier pairs of period 168 h as exogenous regressors (dynamic harmonic regression). \\
SARIMAX & s4 & 0.669 & 11.2 & 139 & Public holidays behave like Sundays; a holiday dummy lets the level drop without disturbing the seasonal state. \\
SARIMAX & s5 & 0.769 & 12.7 & 416 & Check whether more harmonics (5) and a second AR lag improve or overfit; AIC vs validation MASE decides. \\
LightGBM & g1 & 0.642 & 10.7 & 7 & Lags 1-24 only: how much can trees recover from the last day alone? Establishes the value of the seasonal/weekly lags added next. \\
LightGBM & g2 & 0.672 & 10.8 & 7 & Add lags 36-168 (two-day and weekly cycle) identified in the ACF. \\
LightGBM & g3 & 0.604 & 9.7 & 10 & Add calendar features of origin and target (hour, weekday, holiday) so the model can anticipate the day type of the target time, which lags alone cannot at 24 h. \\
LightGBM & g4 & 0.664 & 10.6 & 9 & Add rolling means/stds (6, 24, 168 h) and differences: level and volatility summaries that reduce reliance on any single lag. \\
LightGBM & g5 & 0.682 & 10.9 & 6 & Regularise: fewer leaves, lower learning rate and larger leaves to test whether g4 overfits the pooled data (compare train/val gap). \\
LightGBM & g6 & 0.652 & 10.4 & 9 & Opposite direction: more capacity. Together with g5 this brackets the capacity sweet spot before the Optuna search. \\
LightGBM & Optuna (20 trials) & 0.633 & 10.2 & 244 & best trial opt05: num\_leaves=67, learning\_rate=0.0182, min\_child\_samples=32, feature\_fraction=0.796, lambda\_l2=0.00153 \\
RNN & l1 & 0.721 & 11.5 & 12 & Small network with one day of context, to see the floor of the architecture and whether training is stable. \\
RNN & l2 & 0.747 & 12.5 & 6 & One week of context: the weekly cycle should now be visible to the network. \\
RNN & l3 & 0.714 & 11.5 & 6 & Double the hidden size to test whether l2 was capacity-limited. \\
RNN & l4 & 0.793 & 12.6 & 7 & Deeper network with dropout: more capacity, but regularised against the small number of training weeks. \\
RNN & l5 & 0.711 & 11.6 & 7 & GRU has fewer parameters than LSTM at equal hidden size; test whether the gating simplification helps with ~6 weeks of data. \\
RNN & l6 & 0.720 & 11.4 & 12 & Lower learning rate and light weight decay to smooth the validation curve of the best previous configuration. \\
\bottomrule
\end{tabular}
}{(pending)}
\end{table*}

% ------------------------------------------------------------------------------------
\section{Results and Discussion}\label{sec:results}
\subsection{Tuning trail}
Table~\ref{tab:log} summarises the tuning rounds. For SARIMAX, s2 reduced validation MASE from 0.717 to 0.659. Weekly Fourier terms (s3, 0.666), the holiday dummy (s4, 0.669), and the larger s5 model (0.769) did not improve that score, so s2 was selected. For LightGBM, adding target-time calendar features produced the best manual run (g3, 0.604); rolling features worsened validation MASE to 0.664. The best of \optunaTrials{} Optuna trials reached \optunaBestMase{}, which also failed to beat g3. In the recurrent search, the 168-hour LSTM did not improve on the 24-hour version (0.747 versus 0.721). A larger LSTM recovered to 0.714, and the GRU variant l5 was narrowly best at 0.711. These results support the selected configurations, but they also show that several plausible changes based on the time-series analysis did not generalise to the validation week.

\subsection{Comparative performance}
Table~\ref{tab:results} reports test-set accuracy averaged over cells; Table~\ref{tab:details} gives the cell-level breakdown, significance counts and computational cost; Fig.~\ref{fig:mase} visualises the accuracy results.

\begin{table*}[t]
\centering\caption{Test-set accuracy (22 Dec -- 1 Jan), mean over the selected cells. MASE $<1$ beats the 24-hour seasonal naive forecast.}\label{tab:results}
\footnotesize
\IfFileExists{generated/results_table.tex}{\begin{tabular}{l rr rr rr rrr}
\toprule
 & \multicolumn{2}{c}{$h=1$ h} & \multicolumn{2}{c}{$h=6$ h} & \multicolumn{2}{c}{$h=24$ h} & \multicolumn{3}{c}{All horizons} \\
\cmidrule(lr){2-3}\cmidrule(lr){4-5}\cmidrule(lr){6-7}\cmidrule(lr){8-10}
Model & MASE & sMAPE & MASE & sMAPE & MASE & sMAPE & MASE & sMAPE & MAE \\
\midrule
Naive & 0.471 & 13.3 & 2.094 & 54.5 & 0.873 & 24.2 & 1.146 & 30.7 & 1,068.9 \\
Seasonal naive (24 h) & 0.854 & 22.3 & 0.863 & 23.0 & 0.873 & 24.2 & 0.863 & 23.2 & 768.8 \\
Seasonal naive (168 h) & 2.380 & 51.0 & 2.369 & 51.0 & 2.326 & 51.5 & 2.358 & 51.2 & 1,816.0 \\
SARIMAX & 0.399 & 11.0 & 1.542 & 35.9 & 2.341 & 54.5 & 1.427 & 33.8 & 1,058.3 \\
LightGBM & 0.949 & 26.1 & 1.780 & 43.0 & 1.568 & 40.6 & 1.432 & 36.6 & 1,087.2 \\
GRU & 1.012 & 26.8 & 2.699 & 60.0 & 2.957 & 62.9 & 2.223 & 49.9 & 1,610.5 \\
\bottomrule
\end{tabular}
}{(pending)}
\end{table*}
\begin{table*}[t]
\centering
\caption{Detailed test results. Left: MASE by cell and horizon (best model in bold; SN-24 / SN-168 are seasonal-naive periods). Right: Diebold--Mariano tests against SN-24 ($p<0.05$), followed by final-run CPU cost.}\label{tab:details}
\begin{minipage}[t]{0.56\textwidth}
\centering\scriptsize\setlength{\tabcolsep}{3.5pt}
\textbf{(a) Cell-level MASE}\par\smallskip
\IfFileExists{generated/mase_by_cell_table.tex}{\begin{tabular}{l r rrrrrr}
\toprule
Cell & $h$ & Naive & SN-24 & SN-168 & SARIMAX & LightGBM & GRU \\
\midrule
bocconi & 1 & \textbf{0.092} & 0.260 & 1.207 & 0.105 & 0.483 & 0.497 \\
bocconi & 6 & 0.325 & \textbf{0.271} & 1.212 & 0.516 & 1.084 & 1.803 \\
bocconi & 24 & \textbf{0.255} & 0.255 & 1.203 & 1.119 & 0.805 & 2.337 \\
business & 1 & 0.205 & 0.462 & 2.438 & \textbf{0.179} & 0.970 & 0.967 \\
business & 6 & 0.829 & \textbf{0.460} & 2.434 & 0.905 & 1.770 & 2.816 \\
business & 24 & \textbf{0.426} & 0.426 & 2.401 & 2.136 & 1.570 & 3.367 \\
business-2 & 1 & 0.835 & 1.361 & 1.930 & \textbf{0.412} & 0.631 & 0.607 \\
business-2 & 6 & 4.196 & 1.365 & 1.928 & \textbf{1.243} & 1.402 & 1.662 \\
business-2 & 24 & 1.464 & 1.464 & 1.876 & 1.646 & \textbf{1.329} & 1.722 \\
duomo & 1 & 0.737 & 1.378 & 2.008 & \textbf{0.410} & 0.764 & 0.665 \\
duomo & 6 & 3.534 & 1.424 & 1.977 & \textbf{1.323} & 1.447 & 1.731 \\
duomo & 24 & 1.568 & 1.568 & 2.018 & 1.712 & \textbf{1.523} & 1.710 \\
mixed suburban & 1 & 0.539 & 0.896 & 2.857 & \textbf{0.487} & 0.925 & 1.068 \\
mixed suburban & 6 & 2.424 & \textbf{0.898} & 2.855 & 1.966 & 1.959 & 2.693 \\
mixed suburban & 24 & \textbf{0.927} & 0.927 & 2.878 & 3.202 & 1.654 & 3.384 \\
mixed suburban-2 & 1 & \textbf{0.457} & 0.901 & 3.482 & 0.634 & 1.524 & 1.543 \\
mixed suburban-2 & 6 & 1.661 & \textbf{0.875} & 3.440 & 2.558 & 2.391 & 3.840 \\
mixed suburban-2 & 24 & \textbf{0.765} & 0.765 & 3.287 & 3.508 & 2.178 & 4.039 \\
navigli & 1 & \textbf{0.430} & 0.720 & 2.739 & 0.564 & 1.343 & 1.738 \\
navigli & 6 & 1.692 & \textbf{0.746} & 2.740 & 2.280 & 2.408 & 4.352 \\
navigli & 24 & \textbf{0.706} & 0.706 & 2.619 & 3.062 & 1.918 & 4.137 \\
\bottomrule
\end{tabular}
}{(pending)}
\end{minipage}\hfill
\begin{minipage}[t]{0.40\textwidth}
\centering\footnotesize
\textbf{(b) Diebold--Mariano outcomes}\par\smallskip
\IfFileExists{generated/dm_table.tex}{\begin{tabular}{l rrr}
\toprule
Model & Sig. better & Sig. worse & Comparisons \\
\midrule
GRU & 2 & 15 & 21 \\
LightGBM & 2 & 14 & 21 \\
SARIMAX & 7 & 10 & 21 \\
\bottomrule
\end{tabular}
}{(pending)}
\vspace{1.5em}

\textbf{(c) Computational cost}\par\smallskip
\IfFileExists{generated/cost_table.tex}{\begin{tabular}{l rr}
\toprule
Model & Fit + forecast time (s) & Peak RSS (MB) \\
\midrule
Baselines & 0.1 & 1,048 \\
GRU & 12.0 & 1,710 \\
LightGBM & 9.1 & 1,058 \\
SARIMAX & 59.5 & 1,052 \\
\bottomrule
\end{tabular}
}{(pending)}
\end{minipage}
\end{table*}
\begin{figure}[t]\centering
\includegraphics[width=\columnwidth]{\fig{results_mase_by_horizon_and_cell}}
\caption{Test MASE by horizon (left) and by cell (right). Dotted line: seasonal-naive parity.}\label{fig:mase}
\end{figure}

The 24-hour seasonal naive is the strongest method overall, with mean MASE \maseSeasonalNaiveAll{}. It beats every learned model at 6 and 24 hours and remains below 1 at each horizon. SARIMAX is the best learned model overall (\maseSarimaAll{}), only 0.005 ahead of LightGBM (\maseLightgbmAll{}). Its advantage is concentrated at one hour, where its MASE is \maseSarimaHOne{}; LightGBM is the stronger learned model at 24 hours (\maseLightgbmHDay{} versus \maseSarimaHDay{}). The GRU is weakest overall (\maseGruAll{}). The Diebold--Mariano counts lead to the same conclusion: relative to the 24-hour seasonal naive, SARIMAX is significantly better in 7 of 21 cell-horizon comparisons and worse in 10, while LightGBM is better in 2 and worse in 14; the GRU is better in 2 and worse in 15.

The cell-level table shows some variation in the ranking of the learned models, but the baselines win most rows. SARIMAX is usually the strongest learned method at 1 and 6 hours, including the business cells, while LightGBM often becomes the strongest learned method at 24 hours. This supports a horizon effect. It does not establish that the inferred activity-profile label determines the winning model, because the experiment contains only seven cells and the profile labels are descriptive. The result is consistent with the warning from Makridakis et al.\ \cite{makridakis2018}: on short series with regular seasonality, simple benchmarks can be difficult to beat. It also differs from spatial deep-learning studies on the full grid \cite{wang2017,zhang2018}, whose models use far more series and explicit spatial information.

\subsection{Failure analysis}
Fig.~\ref{fig:dailyerr} plots daily mean absolute error across the test period, while Fig.~\ref{fig:resid} shows the one-hour residual distribution. Errors rise around Christmas and New Year, although the peak day differs by model: LightGBM and the seasonal naive peak on 25 December, whereas SARIMAX and the GRU peak on 1 January. The residual means are positive for all three learned models, indicating over-forecasting on average during the holiday test window; the mean biases across cells are about 174 for SARIMAX, 559 for LightGBM and 540 for the GRU. The largest relative errors are GRU forecasts for the Bocconi cell on 23 and 30 December at the 24-hour horizon, where the model predicts ordinary daytime levels despite much lower observed activity. These failures fit the data limitation: the training period contains too few comparable holiday closures to estimate the size of the drop.

\begin{figure}[t]\centering
\includegraphics[width=\columnwidth]{\fig{failure_daily_error}}
\caption{Daily mean absolute error per model over the test period.}\label{fig:dailyerr}
\end{figure}
\begin{figure*}[t]\centering
\includegraphics[width=\textwidth]{\fig{failure_residuals_h1}}
\caption{One-hour-ahead residuals against the actual level.}\label{fig:resid}
\end{figure*}

\clearpage

% ------------------------------------------------------------------------------------
\section{Conclusion and Future Work}
This study built a memory-bounded pipeline for the 20.8\,GB Milan telecom dataset, characterised its temporal and spatial structure, and compared three forecasting families on seven cells. The main empirical result is that the 24-hour seasonal naive remains the best method overall. SARIMAX is the best learned model by a narrow margin and performs well at one hour, while LightGBM is the best learned method at 24 hours. The GRU does not justify its additional complexity on six weeks of fitting data. Model rankings vary across horizons and some cells, but this small, descriptively labelled sample does not show that activity profile alone decides which model wins. Holiday-period over-forecasting is the most consistent failure.

\textbf{Limitations.} Only \nSelected{} cells and one activity were modelled. Two months of data contain few holidays, the validation week contains none, and the test period is unusually difficult. Cluster labels are interpretations of profile shape rather than ground-truth land use. The four-cluster solution was retained for interpretability even though the silhouette criterion preferred two clusters. One final SARIMAX fit, for the mixed-suburban cell, did not converge, so its aggregate test result should be interpreted with caution.

\textbf{Future work.} Train the tree and recurrent models on all 10\,000 cells with cell-type embeddings, which would give the sequence model more training examples. Add spatial neighbours' lags to separate the value of spatial information from the effect of training scale in \cite{wang2017,zhang2018}. A longer dataset with several holiday periods would support a fairer holiday evaluation. Probabilistic forecasts would also match capacity planning better than point forecasts because operators care about the upper tail of demand.

\textbf{Reproducibility.} The complete source code, configuration, executed notebook, experiment log and instructions are available in the GitHub repository \cite{repo}.

\section*{Use of AI Tools}
An AI coding assistant (Cursor) was used to scaffold the repository, draft code for the ingestion pipeline, models and plotting, and to draft the structure of this report. All experimental decisions -- dataset choice, cell selection, model specifications, tuning sequence, and the interpretation of results -- were made and verified by the author, who ran the experiments and edited the text. The author takes full responsibility for the content.

\bibliographystyle{IEEEtran}
\bibliography{references}

\end{document}
